%%%% necessary customization for the appendix and headers
\appendix
\makeatletter
\addtocontents{toc}{\protect\renewcommand\protect\cftchappresnum{\@chapapp\ }}
\addtocontents{toc}{\protect\setlength{\protect\cftchapnumwidth}{\widthof{\textbf{Appendix~B~}}}}
\makeatother
\renewcommand{\thechapter}{\Alph{chapter}}

\chapter{Supplementary Derivations} \label{chap:appendix-a}

This appendix collects short algebraic derivations and supplementary numerical tables that support the main formulas and empirical summaries in Chapters~\ref{chap:foundation_llm},~\ref{chap:foundation_power}, and~\ref{chap:protein_results}. These materials are placed here so that the core chapters can emphasize interpretation and study-design implications without losing mathematical transparency.

\section{Residual-Correction Identity Behind Prediction-Powered Estimation}

Chapter~\ref{chap:foundation_llm} uses the efficient-influence-function perspective for mean estimation with a surrogate label. Let $R$ indicate whether the gold-standard outcome $Y$ is observed, with \(\Pr(R=1)=\pi\), and assume that $R$ is independent of $(Y,\widehat{Y})$. Define \(\mu(\widehat{Y})=\mathbb{E}[Y\mid \widehat{Y}]\). Then \(\theta=\mathbb{E}[Y]=\mathbb{E}\{\mu(\widehat{Y})\}+\mathbb{E}\{Y-\mu(\widehat{Y})\}\).
Because $R$ is independent of $(Y,\widehat{Y})$ and $\mathbb{E}[R]=\pi$,
\[
\mathbb{E}\left[
\frac{R}{\pi}\{Y-\mu(\widehat{Y})\}
\middle|\,
Y,\widehat{Y}
\right]
=
Y-\mu(\widehat{Y}),
\]
so averaging over $(Y,\widehat{Y})$ gives
\[
\theta
=
\mathbb{E}\{\mu(\widehat{Y})\}
+
\mathbb{E}\left[
\frac{R}{\pi}\{Y-\mu(\widehat{Y})\}
\right].
\]
This is exactly the decomposition used by the efficient influence function
\[
\phi(O)
=
\mu(\widehat{Y})-\theta
+
\frac{R}{\pi}\{Y-\mu(\widehat{Y})\}.
\]
Replacing expectations by sample averages yields the generic residual-correction estimator
\[
\widehat{\theta}(\mu)
=
\frac{1}{N}\sum_{i\in U}\mu(\widehat{Y}_i)
+
\frac{1}{n}\sum_{i\in L}\{Y_i-\mu(\widehat{Y}_i)\}.
\]
Standard PPI corresponds to the working approximation $\mu(\widehat{Y})\approx \widehat{Y}$, while linear \texttt{PPI++} corresponds to $\mu(\widehat{Y})\approx \lambda \widehat{Y}$.

\section{Rogan--Gladen Algebra for the Binary Misclassification Model}

Suppose $Y,\widehat{Y}\in\{0,1\}$ and define \(\theta = \Pr(Y=1)\), \(q_0=\Pr(\widehat{Y}=0\mid Y=0)\), and \(q_1=\Pr(\widehat{Y}=1\mid Y=1)\), along with the observed surrogate prevalence \(p=\Pr(\widehat{Y}=1)\).
Using the law of total probability,
\begin{align*}
p
&=
\Pr(\widehat{Y}=1\mid Y=0)\Pr(Y=0)
+
\Pr(\widehat{Y}=1\mid Y=1)\Pr(Y=1) \\
&=
(1-q_0)(1-\theta)+q_1\theta \\
&=
(1-q_0)+\theta(q_0+q_1-1).
\end{align*}
Rearranging gives
\[
\theta(q_0+q_1-1)=p+q_0-1,
\]
and therefore
\[
\theta=\frac{p+q_0-1}{q_0+q_1-1}.
\]
This is the Rogan--Gladen identity used in Chapter~\ref{chap:foundation_llm}. Its stability depends directly on the denominator $q_0+q_1-1$, which explains why the estimator becomes unstable when the surrogate label is only weakly informative.

\section{Mean \texorpdfstring{\texttt{PPI++}}{PPI++} Variance and the Oracle \texorpdfstring{$\lambda$}{lambda}}

Chapter~\ref{chap:foundation_power} uses the one-sample mean estimator
\[
\widehat{\theta}_{\lambda}
=
\overline{Y}_{L}
+
\lambda(\overline{f}_{U}-\overline{f}_{L}),
\]
where $L$ is the labeled sample of size $n$ and $U$ is the prediction-only sample of size $N$. Assuming the labeled and prediction-only samples are independent,
\begin{align*}
\Var(\widehat{\theta}_{\lambda})
&=
\Var(\overline{Y}_{L})
+
\lambda^2\Var(\overline{f}_{U})
+
\lambda^2\Var(\overline{f}_{L})
-2\lambda\Cov(\overline{Y}_{L},\overline{f}_{L}) \\
&=
\frac{\Var(Y)}{n}
+
\lambda^2\frac{\Var(f)}{N}
+
\lambda^2\frac{\Var(f)}{n}
-2\lambda\frac{\Cov(Y,f)}{n}.
\end{align*}
This gives the expanded form
\[
\Var(\widehat{\theta}_{\lambda})
=
\frac{\Var(Y)}{n}
+
\lambda^2\Var(f)\left(\frac{1}{N}+\frac{1}{n}\right)
-2\lambda\frac{\Cov(Y,f)}{n}.
\]
Using
\[
\Var(Y-\lambda f)=\Var(Y)+\lambda^2\Var(f)-2\lambda\Cov(Y,f),
\]
the same variance can be written as
\[
\Var(\widehat{\theta}_{\lambda})
=
\frac{\Var(Y-\lambda f)}{n}
+
\frac{\lambda^2\Var(f)}{N}.
\]

To obtain the oracle tuning parameter, differentiate with respect to $\lambda$:
\[
\frac{\partial}{\partial \lambda}\Var(\widehat{\theta}_{\lambda})
=
2\lambda\Var(f)\left(\frac{1}{N}+\frac{1}{n}\right)
-2\frac{\Cov(Y,f)}{n}.
\]
Setting the derivative equal to zero gives
\[
\lambda^\star
=
\frac{\Cov(Y,f)}
{\Var(f)\left(1+\frac{n}{N}\right)}.
\]
The second derivative is
\[
\frac{\partial^2}{\partial \lambda^2}\Var(\widehat{\theta}_{\lambda})
=
2\Var(f)\left(\frac{1}{N}+\frac{1}{n}\right)>0,
\]
so this critical point is indeed a variance minimizer.

Substituting $\lambda^\star$ into the variance gives
\[
\Var(\widehat{\theta}_{\lambda^\star})
=
\frac{\Var(Y)}{n}
-
\frac{\Cov(Y,f)^2}{\Var(f)}
\frac{N}{n(n+N)}.
\]
If one writes \(\rho_{Yf}=\Corr(Y,f)\) and \(r=\frac{n}{N}\), then the same expression becomes
\[
\Var(\widehat{\theta}_{\lambda^\star})
=
\frac{\Var(Y)}{n}
\left(
1-\frac{\rho_{Yf}^2}{1+r}
\right).
\]
When $N\gg n$, so that $r\approx 0$, this becomes
\[
\Var(\widehat{\theta}_{\lambda^\star})
\approx
\frac{\Var(Y)}{n}(1-\rho_{Yf}^2),
\]
which is the basis of the Chapter~\ref{chap:foundation_power} rule-of-thumb that the labeled-sample reduction behaves like an $R^2$ gain.

\section{Wald Inversion and Required Labeled Sample Size}

For the one-sample mean problem, consider
\[
H_0:\theta=\theta_0
\qquad\text{versus}\qquad
H_1:\theta=\theta_0+\delta.
\]
Define the target variance scale \(S^2=\left(\frac{\delta}{z_{1-\alpha/2}+z_{1-\beta}}\right)^2\).
The planning rule is to choose the smallest labeled sample size such that
\[
\Var(\widehat{\theta}_{\lambda})\le S^2.
\]

\subsection{Fixed-\texorpdfstring{$\lambda$}{lambda} inversion}

If $\lambda$ is treated as fixed, then
\[
\frac{\Var(Y-\lambda f)}{n}
+
\frac{\lambda^2\Var(f)}{N}
\le S^2,
\]
which implies
\[
n
\ge
\frac{\Var(Y-\lambda f)}
{S^2-\lambda^2\Var(f)/N},
\]
provided $S^2>\lambda^2\Var(f)/N$. This is the direct inversion used for vanilla PPI ($\lambda=1$) and for fixed plug-in \texttt{PPI++} using an estimated $\widehat{\lambda}$.

\subsection{Oracle-\texorpdfstring{$\lambda$}{lambda} inversion}

For the oracle case, the tuning parameter depends on $n$ through $r=n/N$. Using the correlation form of the oracle variance,
\[
\frac{\Var(Y)}{n}
\left(
1-\frac{\rho_{Yf}^2}{1+n/N}
\right)
\le S^2.
\]
Multiplying through and rearranging yields the quadratic inequality
\[
S^2 n^2
+
\left(NS^2-\Var(Y)\right)n
-
\Var(Y)N(1-\rho_{Yf}^2)
\ge 0.
\]
Therefore the required oracle labeled sample size is the smallest integer above the positive root
\[
n_{\mathrm{req}}^{\mathrm{oracle}}
=
\frac{
\Var(Y)-NS^2
+
\sqrt{(NS^2-\Var(Y))^2+4S^2\Var(Y)N(1-\rho_{Yf}^2)}
}
{2S^2}.
\]
This makes the design logic explicit: stronger predictions, represented by larger $\rho_{Yf}^2$, shift the required labeled sample size downward.

\section{Dependence-Sensitive Variance Decomposition and Sample-Size Inflation}
\label{sec:appendix_deff}

Chapter~\ref{chap:protein_results} keeps the published i.i.d.\ planning formulas as the baseline and then introduces a design-effect sensitivity analysis for shared-protein dependence. This section records the short algebra behind that sensitivity rule.

For a fixed tuning parameter \(\lambda\), write
\[
\widehat{\theta}_{\lambda}
=
\bar{R}^{\lambda}_{L}+\bar{G}^{\lambda}_{U},
\qquad
R_e^{\lambda}=Y_e-\lambda f_e,
\qquad
G_e^{\lambda}=\lambda f_e,
\]
where \(L\) is the labeled edge set with size \(n\) and \(U\) is the prediction-only edge set with size \(N\). Because \(\bar{R}^{\lambda}_{L}=n^{-1}\sum_{e\in L}R_e^{\lambda}\) and \(\bar{G}^{\lambda}_{U}=N^{-1}\sum_{e\in U}G_e^{\lambda}\),
\[
\Var(\bar{R}^{\lambda}_{L})
=
\frac{1}{n^2}\sum_{e,e' \in L}\Cov(R_e^{\lambda}, R_{e'}^{\lambda}),
\qquad
\Var(\bar{G}^{\lambda}_{U})
=
\frac{1}{N^2}\sum_{e,e' \in U}\Cov(G_e^{\lambda}, G_{e'}^{\lambda}),
\]
and
\[
\Cov(\bar{R}^{\lambda}_{L},\bar{G}^{\lambda}_{U})
=
\frac{1}{nN}\sum_{e \in L}\sum_{e' \in U}\Cov(R_e^{\lambda}, G_{e'}^{\lambda}).
\]
Therefore
\[
\Var_{\mathrm{dep}}(\widehat{\theta}_{\lambda})
=
\Var(\bar{R}^{\lambda}_{L})
+
\Var(\bar{G}^{\lambda}_{U})
+
2\Cov(\bar{R}^{\lambda}_{L}, \bar{G}^{\lambda}_{U}).
\]

Under the i.i.d.\ approximation, the off-diagonal covariance terms vanish and one recovers
\[
\Var_{\mathrm{iid}}(\widehat{\theta}_{\lambda})
=
\frac{\Var(Y-\lambda f)}{n}
+
\frac{\lambda^2 \Var(f)}{N}.
\]
Let
\[
D_{\lambda}
=
\Var_{\mathrm{dep}}(\widehat{\theta}_{\lambda})
-
\left\{
\frac{\Var(Y-\lambda f)}{n}
+
\frac{\lambda^2 \Var(f)}{N}
\right\}
\]
denote the excess variance beyond that baseline. If \(\Var(Y-\lambda f)>0\) and \(D_{\lambda}\ge 0\), define
\[
d_{\lambda}
=
1+\frac{nD_{\lambda}}{\Var(Y-\lambda f)}.
\]
Substituting
\[
D_{\lambda}
=
\frac{(d_{\lambda}-1)\Var(Y-\lambda f)}{n}
\]
into the dependence-aware variance gives the compact representation
\[
\Var_{\mathrm{dep}}(\widehat{\theta}_{\lambda})
=
d_{\lambda}\frac{\Var(Y-\lambda f)}{n}
+
\frac{\lambda^2 \Var(f)}{N}.
\]

Now define
\[
S^2
=
\left(
\frac{\Delta}{z_{1-\alpha/2}+z_{1-\beta}}
\right)^2
\]
for the usual two-sided Wald design targeting absolute difference \(\Delta\) at significance level \(\alpha\) and power \(1-\beta\). The planning rule is to choose the smallest labeled sample size such that
\[
\Var_{\mathrm{dep}}(\widehat{\theta}_{\lambda})\le S^2.
\]
If \(S^2>\lambda^2\Var(f)/N\), then solving this inequality for \(n\) yields
\[
n_{\mathrm{req}}^{\mathrm{dep}}(\lambda;d_{\lambda})
=
\frac{
d_{\lambda}\Var(Y-\lambda f)
}{
S^2-\lambda^2\Var(f)/N
}.
\]
The corresponding i.i.d.\ target is
\[
n_{\mathrm{req}}^{\mathrm{iid}}(\lambda)
=
\frac{
\Var(Y-\lambda f)
}{
S^2-\lambda^2\Var(f)/N
},
\]
so the dependence-sensitive requirement is simply
\[
n_{\mathrm{req}}^{\mathrm{dep}}(\lambda;d_{\lambda})
=
d_{\lambda}\,n_{\mathrm{req}}^{\mathrm{iid}}(\lambda).
\]
This is the linear sample-size inflation rule used in the Chapter~\ref{chap:protein_results} sensitivity curves. If an integer sample size is required in practice, one takes the ceiling of the real-valued expression above.

\section{Supplementary Numerical Tables for Chapter~\ref{chap:protein_results}}

The appendix now shifts from supporting algebra to supporting numerical detail. The compact Chapter~\ref{chap:protein_results} tables emphasize the main point estimates needed for the narrative; the full interval estimates for those same held-out inference summaries are listed below so that the main chapter can remain interpretive rather than encyclopedic.

\MaybeInput{generated_tables/appendix_chap5_ppi_core_full.tex}

\MaybeInput{generated_tables/appendix_chap5_confusion_full.tex}
